% 09-conclusion.tex 中文对照版（译自 sections/09-conclusion.tex；不加 \label）
\section{结论}

我们描述了 \system{}，一个把 prompt 拼装当作控制面（control surface）的认知层。其贡献在于工件的形态，而不是新机制：一份可执行且会大声失败的契约（14 个文件中 892 条 Python 断言、396 条插件侧断言），一套“恰好一个组件是 fail-closed”的失败语义，以及一种“每个参数的 provenance 都明确、评测仪器被冻结”的评测纪律。机制本身取自既有工作——ACT-R base-level 激活、STPA 不安全控制动作、Ashby 的必要多样性、PID 式抗饱和——我们把它们作为改造（adaptations）呈现，而不是发明。

我们呈现的证据是结构性的，而不是统计性的。契约会执行、并会大声失败；恰好一个节点能够中止一轮，它发出的每个拦截都可归因到一条具名的不安全控制动作；降级会被记录，而不是静默发生。在可以做结构比较的地方，控制面改变了组装边界上的结果：被取代的事实被抑制而不是持续累积；一次静态规则补丁把二十条变异载荷的拦截从 3/20 提升到 20/20。同一套仪器也界定了主张的边界：留出改写仍为 0/40，因此该闸门是词法与结构层面的控制，而不是语义保证。在真实流量规模上确立效应量，是下一项测量；冻结的仪器与开源的回放工具之所以存在，是为了让这项测量成为可能，而不是替代它。
